\setcounter{section}{1}
\section{Constitucion Española de 1978 II}

Nombre completo:  La Constitución Española de 1978 (II). El Poder Judicial. La
justicia en la Constitución. El Consejo General del Poder Judicial.
El Ministerio Fiscal. El Tribunal Constitucional. Naturaleza,
organización y atribuciones.

Apuntes de ASTIC 2020

\localtableofcontents

\subsection{El poder Judicial} % 2
    Es el conjunto de órgamnos que tienen atribuida la función jurisdiccional en la constitución y las leyes.

    Tiene una composición diferenciada respecto a otros poderes del estado. Esto es mas difuso, puesto que es un conjunto de órganos muy diversos y complejos.

    La Constitución otorga la Función Jurisdiccional a Jueces y Magistrados, esta función se entiende como la de juzgar y hacer ejecutar lo juzgado, es decir, resolución de conflictos sociales. Ejerce el control de la legalidad y complementa el ordenamiento jurídico mediante jurisprudencia.

    \subsubsection{La administración de justicia} % 3
        Emana del pueblo, y se administra en nombre del Rey por Jueces y Magistrados integrantes del poder judicial, independientes, inamovibles, responsables y sometidos únicamente al imperio de la ley.

        Es exclusivamente a dichos jueces y magistrados a quien corresponde el ejercicio de la potestad Jurisdiccional.

    \subsubsection{El consejo general del poder judicial} % 3
        El estado se organiza territorialmente en municipios, partidos, provincias y comunidades autónomas, sobre los que ejercen potestad jurisdiccional los Juzgados:
        \begin{itemize}
            \item de Paz
            \item de Primera Instancia e Instrucción
            \item de lo Contencioso Administrativo
            \item de lo Social, de Vigilancia Penitenciaria y de Menores
            \item Audiencias Provinciales
            \item Tribunales Superiores de Justicia
        \end{itemize}
        Y en territorio nacional, la Audiencia Nacional y el tribunal Supremo (Art. 123 de la CE)

\subsection{La justicia en la constitución} % 4
    El Art. 117 da inicio al título donde se definen los principios básicos que regulan a jueces y magistradosL
    \begin{description}
        \item[Principio de independencia] de jueces y magistrados en el ejercicio de la potestad jurisdiccional
        \item[Principio de inamovilidad] No podrán ser separados,suspendidos, trasladados ni jubilados sino por causas y garantías previstas en la ley
        \item[Principio de responsabilidad] por hechos cometidos en el ejercicio de sus funciones, y subsidiaria en que puede incurrir el Estado
        \item[Principio de sometimiento a la ley]
        \item[Principio de unidad jurisdiccional] 
    \end{description}

    También se tienen en cuenta otros principios:
    \begin{description}
        \item[Gratuidad de la Justicia]
        \item[Publicidad y oralidad de las actuaciones judiciales]
        \item[Participación popular]
        \item[Obligación de colaboración]
        \item[Indemnización por error judicial]
        \item[Policía judicial]
        \item[límites en las actuaciones extrajudiciales de jueces y magistrados]
    \end{description}

\subsection{El consejo general del poder judicial} % 6
    Es el órgano de gobierno del poder judicial. Se establece su estatuto y régimen or ley orgánica. Se preside por el presidente del tribunal supremo, y veinte miembros por periodo de 5 años. Doce jueces y magistrados de todas las categorías, 4 por congreso y 4 por senado por mayoría de tres quintos entre abogados y jurístas.

    El CGPJ no es un tribunal, no dicta sentencias ni hace ejecutar. Es un órgano de gobierno interno de organización y administración.

    \subsubsection{Funciones} % 7
        El CGPJ realiza actos administrativos referentes al ingreso, formación, promoción, ascenso y disciplina de jueces y magistrados. Así como la inspección de jueces y tribunales.

        Esto se detalla en el ARt. 560 de la LOPJ y se clasifican en funciones de nombramiento, interposición de conflictos, función de gobierno y de elaboración de normativa. También elabora y ejecuta su propio presupuesto y recibe quejas de los ciudadanos

    \subsubsection{Organización} % 7
        Se establece que el CGPJ se compone de los siguientes óírganos:
        \begin{itemize}
            \item Presidente
            \item Pleno
            \item Comisión permanente
            \item Comisión Disciplinaria y Promotor de la Acción Disciplinaria
            \item Asuntos Económicos
            \item Comisión de Igualdad
            \item Órganos Técnicos
            \begin{itemize}
                \item Secretaría General
                \item Servicio de Inspección
                \item Gabinete Técnico
                \item Escuela Judicial
                \item Centro de Documentación Judicial
                \item Oficina de Comunicación
            \end{itemize}
        \end{itemize}

    \subsubsection{Presidente del consejo general del poder judicial} % 7
        Es necesario ser miembro de la carrera judicial con la categoría de Magistrado del Tribunal Supremo, y reunir las condiciones para ser Presidente de Sala del mismo, o bien ser un jurista de reconocida competencia con mas de 25 años de antigüedad en la profesión.

        En la sesión constitutiva (que será presidida por el vocal de mas edad) deberán presentarse las candidaturas, sin que cada vocal pueda presentar mas de un nombre.

        La elección tendrá lugar en una sesión a celebrar entre 3 y 7 díás después, siendo elegido por mayoría de tres quintos, si no lo es, se votará entre los dos candidatos mas votados a mayoría símple.

        El presidente será nombrado por el Rey mediante Real Decreto, jurará ante el Rey y tomará posesión del cargo ante el pleno del tribunal.

        El mandato coincide con el consejo que lo ha elegido, pudiendo ser reelegido una sola vez.

        Si procede de la carrera judicial, queda en situación de servicios especiales. De lo contrario, será la que corresponda a su cuerpo de procedencia.

        Sus funciones son:
        \begin{itemize}
            \item Convocar y presidir la Sala de Gobierno
            \item Fijar el órden del día de la Sala de Gobierno
            \item Someter propuestas de cometencia a la Sala de Gobierno
            \item Autorizar los acuerdos de la Sala de gobierno y velar por el cumplimiento
            \item Despachar los informes que le pida el CGPJ
            \item Adoptar medidas necesarias ante situaciones de urgencia, dando cuenta en la Sala de Gobierno
            \item Dirigir la inspección de juzgados y Tribunales
            \item Determinar el reparto de asuntos entre salas del Tribunal y sus secciones
            \item Presidir la reunión de presidentes de Salas y Magistrados
            \item Ejercer los poderes dirigidos al orden del Tribunal o Audiencia
            \item Comunicar al CGPJ las vacantes judiciales y de personal auxiliar
            \item Oir las quejas que les hagan los interesados en causas pleitos
            \item Otras previstas en la ley
        \end{itemize}

\subsection{El ministerio fiscal} %10
    Se encarga de promover la acción de la justicia en defensa de la legalidad, de los derechos de los ciudadanos y del interés público.. Así como velar por la independencia de los tribunales y procurar ante estos la satisfacción del interés social.

    El ministerio fiscal actúa ante tribunales de justicia, pero no está integrado en ellos. Tiene personalidad jurídica propia, integrado con autonomía funcional en el poder judicial, ejerce su misión por medio de órganos propios.

    Obedece el principio de jerarquía, es decir, debe respetar las órdenes que emanan de sus superiores. Últimamente es el Fiscal General del Estado. Los fiscales deben regirse también por el principio de imparcialidad que se exige a toda la administración pública. En caso de discrepancia, los fiscales pueden oponerse razonadamente a las órdenes e instrucciones que atenten contra los principios de actuación.

    El fiscal general es nombrado por el rey a propuesta del gobierno, oido el CGPJ. Existen fiscales en diferentes ámbitos territoriales y creadas en base a criterios funcionales. Se encuentra un Fiscal Jefe, un teniente fiscal y después los restantes fiscales.

    Algunos órganos del Ministerio Fiscal son:
    \begin{itemize}
        \item Fiscal General del Estado
        \item Consejo Fiscal
        \item Junta de Fiscales de Sala
        \item Junta de Fiscales Superiores de las Comunidades Autónomas
        \item Fiscalía del Tribunal Supremo
        \item Fiscalía de la Audiencia Nacional
        \item Fiscalías Especiales (Fiscalía contra la Corrupción y la Criminalidad Organizada)
        \item Fiscalía del tribunal de Cuentas
        \item Fiscalías de las Comunidades Autónomas
        \item Físcalías provinciales
    \end{itemize}

    El mandato del Fiscal General durará 4 años. Podrá ser cesado por petición propia, incurrir en incompatibilidades o prohibiciones, incapacidad, incumplimiento grave o cese por el gobierno. Sus principales funciones son:
    \begin{itemize}
        \item Jefatura Suerior del ministerio Fiscal
        \item Impartir órdenes internas de la institución
        \item Dirección e inspección del ministerio fiscal
        \item Proponer al gobierno los ascensos y nombramientos
        \item Designar miembros del ministerio fiscal para que intervengan en los órganos jurisdiccionales en los que esté legitimado para intervenir
    \end{itemize}

\subsection{El tribunal constitucional} %14
    Viene regulado por los Arts. 159 a 165. Se encuentra en un solo título con toda su estructura orgánica. Destacando la función como órgano garantía del órden constitucional.

    La definición del Tribunal Constitucional queda reglada a una Ley Orgánica. El tribunal queda a parte del poder Judicial al estar establecido diréctamente en la constitución. 

    Interpreta la constitución. Es un legislador negativo, ya que sus sentencias determinan la inconstitucionalidad de las normas con rango de ley.

    \subsubsection{Organización} %14
        Se compone de 12 miembros:
        \begin{itemize}
            \item cuatro a propuesta del congreso por tres quintos
            \item cuatro a propuesta del senado por tres quintos
            \item dos por el gobierno
            \item dos por el CGPJ
        \end{itemize} 
        
        Deben ser nombrados entre Magistrados, Fiscales, Profesores universitarios, funcionarios públicos y Abogados, todos jurístas de reconocida competencia con mas de quince años de ejercicio profesional.

        Se designan por un periodo de 9 años y se renuevan por tercios cada tres años.

    \subsubsection{Atribuciones} %15
        Tienen jurisdicción en todo el territorio español y son competentes para conocer:
        \begin{itemize}
            \item recurso de inconstitucionalidad
            \item recurso de amparo por violación de derechos y libertados del artículo 53.2 (igualdad)
            \item conflictos de competencia entre Estado y Comunidades Autónomas, o entre ellas
            \item otros
        \end{itemize}

        La inconstitucinalidad se regla por el Art. 163. Donde se comprueba a solicitud de un juez, si la ley a aplicar en un procedimiento, se ajusta a constitución.

    \subsubsection{Legitimidad para interponer recurso de inconstitucionalidad} %16
        Están legitimados para interponer el recurso de inconstitucionalidad:
        \begin{itemize}
            \item Presidente del gobierno
            \item Defensor del pueblo
            \item 50 diputados
            \item 50 senadores
            \item órganos colegiados ejecutivos de las comunidades autónomas y las asambleas
        \end{itemize}

        Para interponer recurso de amparo:
        \begin{itemize}
            \item Toda persona que invoque un interés legítimo
            \item defensor del pueblo
            \item Ministerio fiscal
        \end{itemize}

    \subsubsection{Tutela por el tribunal constitucional} %16
        Cualquier ciudadano podrá recabar la tutela ante tribunales ordinarios, y en su caso a través del recurso de amparo ante el tribunal constitucional.

\subsection{Normativa TIC en el Ámbito de la administración de justicia} %17
    Con la Ley 18/2011 se pretende generalizar el uso de tecnologías para profesionales de justicia y definir los requisitos de interconexión, interoperabilidad y seguridad para el desarrollo de los aplicativos del mundo judicial.

    La administración de justicia presenta características distintas de otras administraciones públicas. La relación con los ciudadanos es mayormente a través de profesionales. Se consideró que la ley 11/2007 no era plenamente aplicable a la justicia y requiere una regulación específica. Esta ley sólo aplica a las normas procesales y no al procedimiento administrativo

    La ley consta de los siguientes títulos:
    \begin{description}
        \item[Título I] definición del objeto de la ley
        \item[Título II] Regulación del uso de medios electrónicos en la administración de justicia. Derechos y libertades de los profesionales en sus relaciones con la administración por medios electrónicos.
        \item[Título III] Régimen jurídico de la administración judicial electrónica. \begin{itemize}
            \item Sedes judiciales electrónicas. A través de estas se realizan actuaciones por parte de ciudadanos y profesionales
            \item formas de identificación y autenticación de ciudadanos, profesionales y administración.
        \end{itemize}
        \item[Título IV] Condiciones para la tramitación íntegra de los procedimientos judiciales de forma electrónica. \begin{itemize}
            \item Expediente judicial electrónico
            \item registro de escritos, comunicaciones y notificaciones electrónicas
            \item tramitación electrónica de procedimientos judiciales
        \end{itemize}
        \item[Título V] Aspectos básicos sobre la cooperación y colegiación de esfuerzos entre administraciones.
    \end{description}

    \subsubsection{Expediente judicial electrónico } %18
        Es el conjunto de documentos electrónicos correspondientes a un procedimiento judicial. Se asigna un número de identificación general a aquellos documentos que puedan generar un nuevo procedimiento. Este es único e inalterable, permitiendo su identificación unívoca en todo el proceso y entre administraciones.

        El foliado de expedientes se llevará a cabo mediante un índice electrónico. Garantizará la integridad del expediente y permitirá su recuperación. Es admisible que un documento forme parte de varios expedientes distintos.

        La remisión de expedientes se sustituye por la puesta a disposición del expediente judicial electrónico.

    \subsubsection{Esquema Judicial Interoperable y Seguro} %18
        Será aplicado en la administración de justicia para asegurar el acceso, integridad, disponibilidad, autenticidad, confidencialidad, trazabilidad y conservación de los datos.

        Se persigue facilitar las TIC en la actividad administrativa judicial de forma eficaz, eficiente, racional, y de calidad. También permitir a los juzgados y tribunales operar entre sí y con el Ministerio Público.

    \subsubsection{INSIDE} %20
        Servicio desarrollado por el Ministerio de Asuntos Económicos y Transformación Digital. Permite a todos los usuarios una vía de remisión de expedientes administrativos electrónicos a la Administración de Justicia

    \subsubsection{La Oficina Judicial} %20
        Organización instrumental de soporte y apoyo a la actividad jurisdiccional de jueces y tribunales conforme a la LO 19/2003 PJ

        Asume las siguientes actividades
        \begin{description}
            \item[jurisdiccional] por jueces y magistrados
            \item[procedimental] letrados de la administración de justicia y servicios de apoyo y procesales
            \item[administrativa] por el Ministerio de Justicia o Comunidades Autónomas con competencias asumidas
        \end{description}

    \subsubsection{Sistemas de gestión procesal} %21
        Aplicaciones que soportan la tramitación de información relativa a los procedimientos judiciales. De esta forma, cualquier órgano judicial implicado en un determinado procedimiento, puede acceder a la información asociada con garantías de ley

        Ofrece soporte a la estructura organizativa, incluyendo la posibilidad de configurar servicios comunes procesales.

        Un ejemplo es Avantius, que ha permitido agilizar la justicia en Aragón mediante la implantación del expediente judicial electrónico.

    \subsubsection{LexNET} %21
        Sistema de gestión de notificaciones telemáticas desde los juzgados a los profesionales de justicia. Regulado en el RD 1065/2015 sobre comunicaciones electrónicas en la administración de justicia.

    \subsubsection{Planes de modernización de la justicia} %23
        Ha habido numerosos planes para modernizar la justicia y evitar las imágenes de montañas de papeles.

        \paragraph{Plan de acción de la justicia 2012-2015} %23
            Plan de 5 líneas estratégicas con actuaciones y proyectos agrupados en programas.
            
            Destaca el programa 2: Infraestructuras tecnológicas y servicios de actividad judicial. Incluye los siguientes proyectos: Evolución y Coordinación del Expediente Judicial Electrónico en la Audiencia Nacional

            El Programa 3 pretende la convergencia entre sistemas y aplicaciones que prestan servicios en la administración de justicia
            \begin{itemize}
                \item Evolución e implantación de LexNET
                \item Cooperación con organismos nacionales mediante conexión y remisión de datos entre el Sistema Integrado de Registros de Apoyo a la Actividad Judicial.
                \item Cooperación con organismos internacionales como el acceso transfronterizo de ciudadanos y empresas a los medios legales de europa o la interconexión de registros europeos.
            \end{itemize}

            El programa 4: Procedimientos de Gestión de la Administración de Justicia, destaca:
            \begin{itemize}
                \item Cuadro de Mando para Inspección y Auditoría
                \item Estrategia Web en la Administración de Justicia
                \item Portal de subastas
            \end{itemize}

        \paragraph{Plan 2015 para aceleración de la justicia en entornos digitales} %23
            Se pretende conseguir una justicia de Calidad, Ágil y al alcance, que se adapte al S. XXI. Tiene como objetivos estratégicosL Maximizar la eficacia y eficienciencia, orientar el servicio al ciudadano y sus demandas y asentar la transformación en la cultura de innovación y gestión.

            Sus ejes de transformación son la justicia Digital, en Red y Abierta.

            La \textbf{Justicia Digital} conlleva una justicia sin papeles. Se obliga a todos los profesionales a usar los sistemas telemáticos existentes para la presentación de escritos y documentos. Se establece un primer emplazamiento electrónico obligatorio para personas jurídicas y se otorga a los ciudadanos el derecho a elegir cómo comunicarse con la justicia, sea o no de manera electrónica.

            La \textbf{Justicia en Red} implica la remisión electrónica desde los juzgados y las bases de interoperabilidad

            La \textbf{Justicia abierta} en cuanto a la relación con los ciudadanos, se caracteriza por la justicia móvil (avisos por sms, y APP lexNET), y la sede digital 24/7

        \paragraph{Plan de justicia 2030} %25
            Iba a ser aprobado, cuando la pandemia lo obligó a sustiturse por un plan de choque. Destacan el acc eso a la Justicia consolidando garantías y derechos, promover la eficiencia operativa del servicio público de la justicia y la transformación digital.

    \subsubsection{Plan E-Justice de la unión europea} %25
        Se ve en el tema 5.

        Destaca el portal e-justice pensado para ser ventanilla única donde se puede encontrar información dirigida a los siguientes actores:
        \begin{itemize}
            \item Ciudadanos \begin{itemize}
                \item Facilitar información para encontrar abogados y notarios
                \item Órgano judicial al que acudir en cada estado miembro, leyes aplicable, costas, asistencia jurídica\dots
            \end{itemize}
            \item Empresas \begin{itemize}
                \item Registros de entidades insolventes
                \item Registros mercantiles
                \item Requerimientos de pago
                \item Registros de propiedad
                \item Información sobre los requerimientos de profesiones jurídicas
            \end{itemize}
            \item Profesionales del derecho \begin{itemize}
                \item información sobre el derecho de la UE, derecho internacioal y derecho de los Estados Miembros
                \item Profesiones jurídicas y redes de justicia
                \item Red de justicia Europea en materias civil y mercantil
            \end{itemize}
            \item Jueces y Magistrados \begin{itemize}
                \item Información general sobre derecho de la UE, internacional, y de los estados miembros
                \item Red de justicia Europea 
                \item Información sobre uso de tecnologías como videoconferencias
            \end{itemize}
        \end{itemize}

        E-Sens son los Servicios electrónicos europeos. Responde a la demanda de mecanismos interoperables para los servicios públicos digitales.
% Resumen % 27